% Prepared for Edward Baker, 22 October 2026 (America/New_York).
% Substantial LLM assistance: gpt-6-astra (Codex), reasoning effort ultra.
% Module for the investigations-4 analytic manuscript. No new macros required.
% Required environments: theorem, proposition, lemma, corollary, remark, proof.
% Citation keys supplied separately in i4_statistics_caps_bibliography.tex.

\section{Pair statistics and the inverse-square observable}
\label{sc:statistics}

A statistical description becomes useful for heat-flow topology only after
its normalization, limiting procedure, and singular weights have been
specified. We give a precise formulation and prove the relevant transfer
statements. All populations in this section are finite, with labels
retained when locations repeat.

Let $x_1,\ldots,x_m\in\mathbb R$, let $\rho>0$, and put
$y_i=\rho x_i$. The factor $\rho$ is the chosen unfolding density. Define
\begin{equation}
 \mu_m=\frac1m\sum_{i\ne j}\delta_{y_i-y_j}.
 \label{sc:pairmeasure}
\end{equation}
The sum is over ordered pairs of distinct labels. Thus a repeated location
contributes an atom at zero even though the label diagonal is omitted.
The sine-kernel pair density is
\begin{equation}
 R_2(s)=1-\left(\frac{\sin\pi s}{\pi s}\right)^2,
 \qquad R_2(0)=0.
 \label{sc:sine}
\end{equation}
A sine-kernel pair limit means vague convergence
$\mu_m\to R_2(s)\,ds$, that is, convergence against every continuous
compactly supported test function. Its meaning includes the choice of
populations and of the density $\rho$.

Montgomery's original formulation concerns a growing height cutoff and
unfolding by $\log T/(2\pi)$, with RH assumed in that paper; $T$ there is
not a heat time. His discussion identifies the sine-kernel pair law and
its random-matrix connection, and obtains density-one simplicity from
the conjectured law \cite[pp.~181--184]{sc:Montgomery1973}. In a heat
problem one must separately prescribe the joint limits of heat time,
height, and population. Fixed-rank late-time information alone does not
specify a growing-population law.

If the points are distinct, define the real antisymmetric matrix
\begin{equation}
 A_{ij}=\begin{cases}(y_i-y_j)^{-1},&i\ne j,\\0,&i=j,\end{cases}
 \qquad
 Q_m=\frac1m\lVert A\rVert_F^2
     =\int_{\mathbb R}\frac{d\mu_m(s)}{s^2}.
 \label{sc:matrixenergy}
\end{equation}
Here the integrand is assigned the value $+\infty$ at zero, so the last
expression also defines an extended nonnegative quantity when labels
repeat. The observable $Q_m$ is the ordered inverse-square pair energy
per point. In physical coordinates,
\begin{equation}
 8\sum_{i<j}\frac1{(x_i-x_j)^2}=4m\rho^2Q_m.
 \label{sc:physicalweight}
\end{equation}
This is the pair term in the logarithmic derivative of the cluster
discriminant. The complete heat formula also contains the exterior
holomorphic factor; changing to matrix notation does not remove that
factor. Moreover, at complex time the holomorphic weight
$(z_i-z_j)^{-2}$ is distinct from the positive weight
$|z_i-z_j|^{-2}$ in a Frobenius norm.

\subsection{Exact clock and sine-kernel moments}
\label{sc:moments}

\begin{proposition}[Two different inverse-square target values]
\label{sc:targets}
For $m$ consecutive unit-spaced points,
\begin{equation}
 Q_m^{\mathrm{clock}}=
 \frac2m\sum_{d=1}^{m-1}\frac{m-d}{d^2}
 \longrightarrow\frac{\pi^2}{3}.
 \label{sc:clockmoment}
\end{equation}
The sine-kernel target instead has the finite moment
\begin{equation}
 \int_{\mathbb R}\frac{R_2(s)}{s^2}\,ds=
 \frac{2\pi^2}{3}.
 \label{sc:sinemoment}
\end{equation}
\end{proposition}

\begin{proof}
There are $m-d$ unordered pairs at distance $d$, which proves the finite
formula in \eqref{sc:clockmoment}. Rewriting it as
\[
 2\sum_{d=1}^{m-1}d^{-2}
 -\frac2m\sum_{d=1}^{m-1}d^{-1}
\]
and using $\sum_{d\ge1}d^{-2}=\pi^2/6$ proves its limit. The latter
identity follows, for example, by applying Parseval to the Fourier sine
series of $x$ on $[-\pi,\pi]$, whose coefficients are
$2(-1)^{d+1}/d$.

For the second claim, the elementary triangular Fourier identity is
\begin{equation}
 \left(\frac{\sin\pi s}{\pi s}\right)^2
 =\int_{-1}^{1}(1-|u|)\cos(2\pi us)\,du.
 \label{sc:triangular}
\end{equation}
It follows either by integrating the right side twice, or by taking the
Fourier transform of the convolution of two interval indicators. Since
$\int_{-1}^{1}(1-|u|)\,du=1$, the integrand in
\[
 \frac{R_2(s)}{s^2}
 =\int_{-1}^{1}(1-|u|)
       \frac{1-\cos(2\pi us)}{s^2}\,du
\]
is nonnegative. Tonelli's theorem therefore applies without a prior
integrability assumption. The elementary integral
\[
 \int_{\mathbb R}\frac{1-\cos(as)}{s^2}\,ds=\pi|a|
\]
may be obtained by scaling and integration by parts, using the
Dirichlet integral $\int_0^\infty\sin v/v\,dv=\pi/2$; the latter follows
by first inserting $e^{-\varepsilon v}$, differentiating with respect
to the sine frequency, and then letting $\varepsilon\downarrow0$.
Consequently
\[
 \int_{\mathbb R}\frac{R_2(s)}{s^2}\,ds
 =2\pi^2\int_{-1}^{1}|u|(1-|u|)\,du
 =\frac{2\pi^2}{3}.
\]
This calculation also proves finiteness. Directly,
$R_2(s)=\pi^2s^2/3+O(s^4)$ at zero and $0\le R_2(s)\le1$ at infinity.
\end{proof}

The normalization difference in Proposition~\ref{sc:targets} is
substantive. A clock and a sine-kernel population have different positive
inverse-square energies; neither target energy is zero. If, for each
fixed finite group of ranks, $L_t x_j(t)/\pi\to j-\tfrac12$ as
$t\to\infty$, then taking the heat limit first and the population limit
second gives \eqref{sc:clockmoment}. A different simultaneous limit
requires estimates uniform in its growing ranks and population.

\subsection{Sparse exceptional points and the information in a pair law}
\label{sc:sparse}

\begin{lemma}[Sparse modifications preserve a local pair limit]
\label{sc:sparselemma}
For each $n$, let $y_1^{(n)},\ldots,y_n^{(n)}$ be a real labeled array.
Suppose its pair measures \eqref{sc:pairmeasure} converge vaguely to a
locally finite measure $\mu$. There is a choice of one old point in each
array such that adjoining a second label at that same location preserves
the vague limit. Alternatively, one can adjoin a distinct point at
distance at most $n^{-2}$ from the chosen point, preserving the vague
limit but ensuring
\begin{equation}
 Q_{n+1}\ge\frac{2n^4}{n+1}.
 \label{sc:nearpairgrowth}
\end{equation}
If the original arrays are simple, the latter arrays can also be chosen
simple.
\end{lemma}

\begin{proof}
For every fixed positive integer $R$, dominate
$\mathbf 1_{[-R,R]}$ by a nonnegative continuous compactly supported
function. Vague convergence gives
\[
 \#\{(i,j):i\ne j,\ |y_i^{(n)}-y_j^{(n)}|\le R\}=O_R(n).
\]
A diagonal selection yields $R_n\to\infty$ sufficiently slowly that
this count with $R=R_n$ is $o(n^2)$. For completeness, choose integers
$N_k$ so large that for $n\ge N_k$ the count at radius $k$ is at most
$n^2/k$, make $N_k$ increasing, and let $R_n$ be the largest admissible
$k$ with $N_k\le n$. Averaging over the first label gives an index $i_n$
with $o(n)$ neighbors inside radius $R_n$.

Adjoin a point $v_n=y_{i_n}^{(n)}$. For a continuous function $f$
supported in $[-R,R]$, the new pair integral equals
\begin{equation}
 \frac{n}{n+1}\int f\,d\mu_n+
 \frac1{n+1}\sum_{j=1}^{n}
 \bigl[f(v_n-y_j^{(n)})+f(y_j^{(n)}-v_n)\bigr].
 \label{sc:addedpoint}
\end{equation}
The first term tends to $\int f\,d\mu$. The sum in the second term has
$o(n)$ nonzero summands, each bounded by $2\lVert f\rVert_\infty$,
so it tends to zero after division by $n+1$. The new repeated pair
contributes the atom $2\delta_0/(n+1)$, which is invisible in the limit.

For the alternative, choose $0<\varepsilon_n\le n^{-2}$ so that
$v_n=y_{i_n}^{(n)}+\varepsilon_n$ is distinct from every old point.
Only finitely many choices are forbidden. For each fixed $R$, its
neighbors within $R$ are old neighbors of $y_{i_n}^{(n)}$ within $R+1$,
so the proof using \eqref{sc:addedpoint} is unchanged. The two ordered
terms associated with $y_{i_n}^{(n)}$ give
\[
 Q_{n+1}\ge\frac{2}{(n+1)\varepsilon_n^2}
 \ge\frac{2n^4}{n+1}.
\]
No existing points were moved, which proves the simplicity assertion.
\end{proof}

Thus weak pair convergence measures the bulk while admitting sparse
exceptional pairs. This is a statement about the information content of
a pair law, independent of whether a particular array can be realized
as a zero set of a specified heat solution.

\begin{corollary}[Density-one simplicity from a sine-kernel pair law]
\label{sc:densityone}
If $\mu_n\to R_2(s)\,ds$, and a distinct location $a$ in the $n$th
array has multiplicity $m_a$, then
\begin{equation}
 \frac1n\sum_a m_a(m_a-1)\longrightarrow0.
 \label{sc:multiplicitydensity}
\end{equation}
In particular, the fraction of labels lying at repeated locations tends
to zero.
\end{corollary}

\begin{proof}
The left side is $\mu_n(\{0\})$. For each $\delta>0$ choose a continuous
function $g_\delta$ with $0\le g_\delta\le1$, $g_\delta(0)=1$, and
support in $[-\delta,\delta]$. Then
\[
 \limsup_{n\to\infty}\mu_n(\{0\})
 \le\int g_\delta(s)R_2(s)\,ds
 \le\int_{-\delta}^{\delta}R_2(s)\,ds.
\]
The final bound tends to zero with $\delta$. At a repeated location
$m_a\le m_a(m_a-1)$, proving the last assertion.
\end{proof}

\subsection{Transfer of the singular inverse-square weight}
\label{sc:transfer}

\begin{proposition}[Singular-moment transfer]
\label{sc:UI}
Suppose $\mu_n\to R_2(s)\,ds$ vaguely and
\begin{align}
 &\lim_{\delta\downarrow0}\limsup_{n\to\infty}
    \int_{|s|<\delta}\frac{d\mu_n(s)}{s^2}=0,
 \label{sc:UIzero}\\
 &\lim_{R\to\infty}\limsup_{n\to\infty}
    \int_{|s|>R}\frac{d\mu_n(s)}{s^2}=0.
 \label{sc:UItail}
\end{align}
Then $Q_n\to2\pi^2/3$.
\end{proposition}

\begin{proof}
Fix $0<\delta<R$. On the compact annulus
$\delta\le|s|\le R$ the weight is bounded, and its only discontinuities
when extended by zero are at four points of zero limiting measure.
Continuous upper and lower cutoffs, or the usual weak-convergence
criterion for bounded functions with a null discontinuity set, give
\[
 \int_{\delta\le|s|\le R}s^{-2}\,d\mu_n
 \longrightarrow
 \int_{\delta\le|s|\le R}s^{-2}R_2(s)\,ds.
\]
The two nonnegative omitted terms are uniformly negligible in the
successive limits by \eqref{sc:UIzero}--\eqref{sc:UItail}. The corresponding
limiting terms vanish by Proposition~\ref{sc:targets}. Removing the
cutoffs proves the assertion.
\end{proof}

Conditions \eqref{sc:UIzero} and \eqref{sc:UItail} identify the exact
additional information needed by this moment transfer. A zero-gap atom
makes the first integral infinite; that condition therefore cannot be
assumed as an innocuous consequence of a weak pair law in an argument
intended to establish simplicity. A regularized measurable observable is
\[
 Q_{n,\varepsilon}=\int_{\mathbb R}
     \frac{d\mu_n(s)}{s^2+\varepsilon^2},\qquad\varepsilon>0,
\]
but its uniform behavior as $\varepsilon\downarrow0$ is an additional
estimate. Conversely, a genuinely proved finite bound $Q_m\le C$ gives
\begin{equation}
 \min_{i\ne j}|y_i-y_j|\ge\sqrt{\frac{2}{mC}},
 \label{sc:mingap}
\end{equation}
because any one pair contributes two terms to $mQ_m$.

\subsection{A finite-population error threshold}
\label{sc:tent}

The following quantitative counting implication specifies how a
sufficiently sharp pair estimate can establish more than density-one
simplicity.

\begin{proposition}[Shrinking-tent test]
\label{sc:tenttheorem}
Let $f_\delta(s)=(1-|s|/\delta)_+$ for $\delta>0$. Suppose a real labeled
population of size $m$ has a proved enclosure
\begin{equation}
 \left|\int f_\delta\,d\mu_m
       -\int_{\mathbb R}f_\delta(s)R_2(s)\,ds\right|
 \le e_m(\delta).
 \label{sc:tenterror}
\end{equation}
If
\begin{equation}
 \frac{\pi^2\delta^3}{18}+e_m(\delta)<\frac2m,
 \label{sc:simplicitythreshold}
\end{equation}
all locations are simple. If the right side in
\eqref{sc:simplicitythreshold} is replaced by $1/m$, every two distinct
labels are separated by more than $\delta/2$.
\end{proposition}

\begin{proof}
By \eqref{sc:triangular} and $1-\cos v\le v^2/2$,
\[
 R_2(s)\le2\pi^2s^2
       \int_{-1}^{1}u^2(1-|u|)\,du
       =\frac{\pi^2s^2}{3}.
\]
Since
\[
 \int_{-\delta}^{\delta}(1-|s|/\delta)s^2\,ds
 =2\left(\frac{\delta^3}{3}-\frac{\delta^3}{4}\right)
 =\frac{\delta^3}{6},
\]
we obtain
\begin{equation}
 \int f_\delta(s)R_2(s)\,ds\le\frac{\pi^2\delta^3}{18}.
 \label{sc:tentconstant}
\end{equation}
Any repeated location contributes at least two ordered terms of value
$f_\delta(0)=1$, hence $\int f_\delta\,d\mu_m\ge2/m$.
Equations \eqref{sc:tenterror}--\eqref{sc:tentconstant} contradict this
under \eqref{sc:simplicitythreshold}. Likewise a pair at distance at
most $\delta/2$ contributes two terms each at least $1/2$, giving the
threshold $1/m$.
\end{proof}

For example, $\delta_m=o(m^{-1/3})$ and
$e_m(\delta_m)=o(m^{-1})$ suffice for eventual simplicity of these finite
populations. This is a conditional quantitative test: neither the
shrinking-test uniformity nor its error bound follows from vague PCC.
The threshold states an analytic task rather than an estimate for the
canonical theta source.

\section{Physical projective caps and their phase}
\label{sc:caps}

We next derive a boundary formula that identifies the information beyond
root positions needed for the physical Chern--Hopf integral. Let
$\tau=t+is$ and $z=x+iy$. Let $H(\tau,z)$ be holomorphic and put
\begin{equation}
 G=(H,H_z),\qquad
 \mathcal A=\frac{\operatorname{Im}(G^*dG)}{|G|^2},\qquad
 \mathcal C=\frac{\mathcal A\wedge d\mathcal A}{4\pi^2}.
 \label{sc:forms}
\end{equation}
These are the physical jet and its normalized connection and three-form;
in particular, the derivative component has not been rescaled relative
to the value component.

\begin{lemma}[Normalization of the boundary three-form]
\label{sc:degreeform}
If $W$ is a compact oriented four-dimensional domain with piecewise
$C^1$ boundary, $G=(H,H_z)$ is nonzero on that boundary, and
$\partial W$ has the outward orientation, then
(with the complex orientation when $W\subset\mathbb C^2$)
\[
 \int_{\partial W}\mathcal C
 =\deg\left(\frac G{|G|}:\partial W\longrightarrow S^3\right).
\]
For a proper product window satisfying Theorem~\ref{tr:complex-charge},
this is the discriminant charge in \eqref{tr:charge-residue}.
\end{lemma}
\begin{proof}
On the unit sphere in $\mathbb C^2$, let
$\alpha=\operatorname{Im}(\overline\zeta_1d\zeta_1+
\overline\zeta_2d\zeta_2)$. At $(1,0)$, an oriented orthonormal tangent
basis is $(\partial_{y_1},\partial_{x_2},\partial_{y_2})$, and
$\alpha\wedge d\alpha$ evaluates to $2$. Unitary invariance and the
transitive action of $U(2)$ therefore give
$\alpha\wedge d\alpha=2\,d\operatorname{Vol}_{S^3}$ everywhere.
The area of the unit three-sphere is $2\pi^2$ (by polar integration of
the four-dimensional Gaussian), so this form divided by $4\pi^2$
has integral one. The physical connection is the pullback of $\alpha$
under $G/|G|$, since differentiating the real norm adds no imaginary
part. The degree formula for a normalized volume form proves the first
claim. Proposition~\ref{tr:degree-agreement} proves the last.
\end{proof}

This identifies the total complete-boundary integral. It does not
identify an individual scalar discriminant edge integral from
Proposition~\ref{tr:complex-caps} with an individual physical
three-dimensional cap; the latter is evaluated next.

Fix real numbers $t_0<T$ and a rectangle
\[
 D=\{\tau:t_0\le\operatorname{Re}\tau\le T,
                 \ |\operatorname{Im}\tau|\le Y\},
 \qquad Y>0,
\]
and a number $\varepsilon>0$. Assume there is a holomorphic map
$Z(\tau,q)$ on a neighborhood of
$D\times\{|q|\le\varepsilon\}$ such that
\begin{equation}
 H_z(\tau,Z(\tau,q))\ne0,\qquad
 \frac{H(\tau,Z(\tau,q))}{H_z(\tau,Z(\tau,q))}=q,
 \label{sc:inversecontract}
\end{equation}
and $q\mapsto Z(\tau,q)$ is one-to-one for each $\tau$.
The inverse derivative is nonzero, because differentiating the second
identity gives $(H/H_z)_z Z_q=1$. Thus
$(\tau,q)\mapsto(\tau,Z(\tau,q))$ is an embedded tube, with the same
orientation as $dt\wedge ds\wedge d\operatorname{Re}q\wedge
 d\operatorname{Im}q$. Indeed the spatial Jacobian is $|Z_q|^2>0$.
Its central branch $Z(\tau,0)$ is a simple root.

\begin{proposition}[Exact phase formula for a physical cap]
\label{sc:capformula}
Under \eqref{sc:inversecontract}, set
\[
 h(\tau,q)=H_z(\tau,Z(\tau,q)),\qquad
 k_\varepsilon=\frac{\varepsilon^2}{1+\varepsilon^2}.
\]
Choose a continuous phase $\phi=\operatorname{Im}\log h$ on the parameter
tube. For the forward real-time cap at $t=b$, oriented by increasing
$s$ followed by the complex $q$ orientation,
\begin{equation}
 I_b=\int_{t=b}\mathcal C
 =\frac{k_\varepsilon}{2\pi}
       \bigl[\phi(b+iY,0)-\phi(b-iY,0)\bigr].
 \label{sc:capexact}
\end{equation}
The projective-circle wall $|q|=\varepsilon$ has identically zero
pullback of $d\mathcal A$ and hence of $\mathcal C$.
\end{proposition}

\begin{proof}
The parameter domain is contractible and $h$ is nonvanishing, so it
admits a logarithm. In these coordinates $G=h(q,1)$, whence direct
substitution in \eqref{sc:forms} gives
\begin{equation}
 \mathcal A=d\phi+a(q),\qquad
 a(q)=\frac{\operatorname{Im}(\overline q\,dq)}{1+|q|^2},\qquad
 d\mathcal A=\omega(q):=
 \frac{2\,du\wedge dv}{(1+u^2+v^2)^2},\quad q=u+iv.
 \label{sc:chartforms}
\end{equation}
Both $a$ and $\omega$ use only the two $q$ coordinates, so
$a\wedge\omega=0$. Therefore
\begin{equation}
 \mathcal C=\frac{d\phi\wedge\omega}{4\pi^2}.
 \label{sc:phaseform}
\end{equation}
On a fixed real-time cap, only $\phi_s\,ds$ in $d\phi$ contributes.
For each $s$, $q\mapsto\phi_s(b+is,q)$ is harmonic: it is the imaginary
part of the $s$ derivative of the holomorphic logarithm of $h$.
Angular averaging on every centered circle gives its central value.
Since $\omega$ has radial density,
\[
 \int_{|q|\le\varepsilon}\phi_s(b+is,q)\,\omega(q)
 =\phi_s(b+is,0)\int_{|q|\le\varepsilon}\omega(q).
\]
The curvature mass is
\[
 \int_{|q|\le\varepsilon}\omega
 =\int_0^{2\pi}\int_0^\varepsilon
       \frac{2r}{(1+r^2)^2}\,dr\,d\theta
 =2\pi k_\varepsilon.
\]
Integrating in $s$ proves \eqref{sc:capexact}. On the wall,
$q=\varepsilon e^{i\theta}$ depends on only one real parameter. Thus
the two-form $\omega$ pulls back to zero, proving the wall assertion.
\end{proof}

Equation \eqref{sc:phaseform} retains the common scalar phase. Replacing
$G=h(q,1)$ by its projective representative $(q,1)$ would give
$a\wedge da=0$ and would remove the very term measured by
\eqref{sc:capexact}. It is therefore essential to carry the phase of
$H_z$ along the continued root. This observation concerns the normalized
jet; arbitrary multiplication of $H$ by a time-dependent scalar is not
asserted to preserve the heat equation.

\subsection{The complete regular tube}
\label{sc:tube}

\begin{proposition}[The imaginary-time faces close the cap identity]
\label{sc:closedtube}
In the setting of Proposition~\ref{sc:capformula}, let $J_+$ and $J_-$
be the outward oriented integrals on $s=Y$ and $s=-Y$, respectively.
Then
\begin{align}
 J_+&=-\frac{k_\varepsilon}{2\pi}
       [\phi(T+iY,0)-\phi(t_0+iY,0)],
 \label{sc:topface}\\
 J_-&=\frac{k_\varepsilon}{2\pi}
       [\phi(T-iY,0)-\phi(t_0-iY,0)],
 \label{sc:bottomface}\\
 I_T-I_{t_0}+J_++J_-&=0.
 \label{sc:tubebalance}
\end{align}
If $H$ and the inverse branch have real type, there is a continuous
choice
\[
 \theta(t,Y)=\operatorname{Im}\log
          \frac{h(t+iY,0)}{h(t,0)}
\]
for which
\begin{equation}
 I_t=\frac{k_\varepsilon}{\pi}\theta(t,Y),\qquad
 J_+=J_-=\frac{k_\varepsilon}{2\pi}
           [\theta(t_0,Y)-\theta(T,Y)].
 \label{sc:realtypebalance}
\end{equation}
\end{proposition}

\begin{proof}
The ambient orientation is $dt\wedge ds\wedge du\wedge dv$.
Contraction with the upper outward normal $\partial_s$ gives
$-dt\wedge du\wedge dv$, whereas the lower outward normal
$-\partial_s$ gives $+dt\wedge du\wedge dv$. Apply the same harmonic
radial averaging as in Proposition~\ref{sc:capformula}, now to
$\phi_t$, and integrate in $t$. This yields
\eqref{sc:topface}--\eqref{sc:bottomface}; substitution of
\eqref{sc:capexact} proves \eqref{sc:tubebalance} by cancellation of
its four endpoint values. Equivalently,
$d\mathcal C=\omega\wedge\omega/(4\pi^2)=0$ on this regular chart,
so Stokes gives the same identity, with the zero wall already paid.

For the last claim, $h(t,0)$ is nonzero real on the real segment, and
therefore has constant sign. The logarithm after division by that sign
can be chosen real on the real segment. Conjugation then gives opposite
imaginary parts at $t+iY$ and $t-iY$. Inserting these two values in the
preceding formulas gives \eqref{sc:realtypebalance}.
\end{proof}

This proposition makes a useful distinction: an exactly vanishing
projective wall is compatible with nonzero real-time caps, because the
imaginary-time faces provide the remaining flux. The assertion requires
a fixed projective radius on each tube. If the radius varies with time,
the wall pulls back $dr\wedge d\theta$ terms and requires a separate
calculation. A tube satisfying \eqref{sc:inversecontract} contains no
joint zero of $(H,H_z)$; attaching it to an event-containing region does
not remove that region's other boundary pieces.

\subsection{Exact heat examples}
\label{sc:controls}

We give two explicit calculations within the backward heat convention
$H_\tau=-H_{zz}$. They isolate the respective roles of root regularity
and physical phase without assuming a statistical model.

\begin{proposition}[A heat event followed by simple clock asymptotics]
\label{sc:clockcontrol}
The entire real-type field
\begin{equation}
 H(\tau,z)=e^\tau\cos z+1
 \label{sc:cosineplusone}
\end{equation}
satisfies the backward heat equation. At $(\tau,z)=(0,\pi)$ it has an
ordinary double spatial zero and local complex degree one for the map
$(H,H_z)$. At every positive real time all its spatial zeros are real
and simple, and its late real zero pattern converges to a clock of
spacing $\pi$.
\end{proposition}

\begin{proof}
Direct differentiation gives $H_\tau=e^\tau\cos z=-H_{zz}$.
At $(0,\pi)$, $H=H_z=0$, $H_{zz}=1$, and the complex Jacobian is
\[
 \frac{\partial(H,H_z)}{\partial(\tau,z)}(0,\pi)
 =\begin{pmatrix}-1&0\\0&1\end{pmatrix}.
\]
Its nonzero determinant makes the zero of this holomorphic map simple,
with local degree one. For $t>0$, the equation is
$\cos z=-e^{-t}\in(-1,0)$. If $z=x+iy$ with $y\ne0$ and $\cos z$
is real, then $\sin x=0$, so $|\cos z|=\cosh y>1$, a contradiction.
Thus all zeros are real, and $\sin z\ne0$ at them, proving simplicity.
The two roots in $(0,2\pi)$ are
\[
 x_\pm(t)=\pi\pm\arccos(e^{-t}).
\]
They tend to $\pi/2$ and $3\pi/2$; translation by multiples of $2\pi$
gives the clock limit. For $t<0$ the same local pair is
$\pi\pm i\operatorname{arcosh}(e^{-t})$, displaying explicitly the
transition through the event.
\end{proof}

The source of \eqref{sc:cosineplusone} is a positive atomic cosine
measure. The example proves that late simplicity and clock relaxation
are compatible with earlier event charge in the general heat class; it
makes no assertion about an event for the canonical theta source.

\begin{proposition}[Identical tracked roots with different cap phases]
\label{sc:phasecontrol}
For $a>0$, the two exact real-type heat fields
\begin{equation}
 H_1(\tau,z)=z,\qquad
 H_2(\tau,z)=e^{a^2\tau}\frac{\sin(az)}a
 \label{sc:twoheatfields}
\end{equation}
have the same simple tracked root $z=0$ at every complex time. For any
$0<\varepsilon<1/a$, their root inverse disks are respectively
\begin{equation}
 Z_1(\tau,q)=q,\qquad
 Z_2(\tau,q)=\frac1a\arctan(aq),
 \label{sc:twoinverses}
\end{equation}
and their physical cap integrals at $t=T$, $-Y\le s\le Y$, are
\begin{equation}
 I_{T,1}=0,\qquad
 I_{T,2}=\frac{\varepsilon^2}{1+\varepsilon^2}
             \frac{a^2Y}{\pi}.
 \label{sc:twocaps}
\end{equation}
\end{proposition}

\begin{proof}
Both heat equations follow by differentiation. Their derivatives at the
tracked root are $1$ and $e^{a^2\tau}$, so that root is simple.
The projective quotients are $q=z$ and $q=\tan(az)/a$.
The principal branch of $\arctan(aq)$ is holomorphic on $|q|<1/a$;
it is one-to-one there because its tangent recovers $aq$, and its
derivative is $(1+a^2q^2)^{-1}\ne0$. Thus both inverse disks satisfy
\eqref{sc:inversecontract}. On the second disk,
\[
 h_2(\tau,q)=\frac{e^{a^2\tau}}{\sqrt{1+a^2q^2}},
\]
with the square root equal to one at $q=0$. The central phases are
therefore $0$ and $a^2s$. Formula~\eqref{sc:capexact} gives
\eqref{sc:twocaps}.
\end{proof}

The last example identifies a precise information deficit of a finite
tracked population: its positions do not specify the physical phase
of the derivative along the root. The second field has exterior roots;
no claim is made about nonuniqueness from an entire global zero set.
Together, the pair-energy results and the cap formula provide a concrete
route for future estimates: control the weighted pair statistic in a
stated limit, retain its complete exterior contribution, and then
estimate the complex-time derivative phase and all faces of the chosen
boundary.
